\section{Discussion}

\subsection{What the experiment establishes}

Within the frozen task matrix, the evidence-bound controller repairs more
controlled workflow failures than execute-once and does so without unsupported
claims. The paired confidence interval excludes zero, exact reruns pass, and
the component ablations behave in the directions implied by their removed
contracts. These observations support the systems claim that explicit
failure feedback, persistent memory, preregistration, and evidence gating can
be made jointly testable.

The two \routea{} rounds add a complementary form of evidence. They show that
the campaign can preserve a favorable development result, reveal an
unfavorable unseen result, and continue through a changed mechanism without
rewriting the prior decision. This is closer to the intended use than a
success-only demonstration. It also prevents the systems result from being
mistaken for evidence that the proposed equation-discovery mechanisms work.

\subsection{Why the gates are separate}

An autonomous workflow can fail at multiple levels. A scientific-method gate
asks whether a candidate improves the substantive task on unseen evidence. A
research-system gate asks whether the controller follows a valid process,
recovers specified failures, reproduces results, and limits claims. A paper
gate asks whether reported claims resolve to evidence and the package can be
rebuilt. Passing one gate does not imply passing the others.

This separation explains the route decision. \routea{} fails its method gate,
so no new sparse-regression contribution is claimed. \routeb{} passes its
internal controlled systems gate, so a systems artifact is built. The paper
package may pass deterministic formatting and reproducibility checks, yet it
still does not authorize submission. Novelty, external validity, and venue
fit are judgments outside the deterministic campaign.

\subsection{Design implications}

First, negative-result lineage should be a first-class data type. A parent
hash and mechanism delta make it possible to distinguish a scientific
successor from an ordinary retry. Second, unseen evaluation should be a
one-way boundary. Once revealed, evidence can inform a new round but cannot
change the current round's registered decision. Third, report generation
should be part of completion. Logs are useful for debugging, but they do not
replace a statement of hypothesis, methods, result, limitations, and next
action.

Fourth, claim quality and execution success should remain separate metrics.
The no-evidence-gate ablation completes many tasks while emitting unsupported
claims. A benchmark that scores only terminal execution would miss that
failure. Finally, local-model fallback should be explicit. A deterministic
fallback can preserve progress under a model timeout, provided that it is
recorded and does not silently assume scientific authority.

\subsection{Relation to benchmark validity}

Kapoor et al. argue that agent evaluations should connect costs, model
choices, and meaningful downstream outcomes \cite{kapoor2024agentsmatter}.
Our controlled matrix follows that recommendation only partially. It exposes
the mechanism and records cost, but the terminal outcomes are workflow
invariants rather than independent discoveries. The experiment is therefore
best understood as a mechanism test and integration benchmark. A stronger
external study should freeze tasks authored by an independent team, include
stochastic policy behavior, measure repair correctness through domain
experiments, and compare against additional agent implementations under a
common compute budget.
